% Part: incompleteness
% Chapter: theories-computability
% Section: consis-with-q

\documentclass[../../../include/open-logic-section]{subfiles}

\begin{document}

\olfileid[pt]{inc}{tcp}{con}

\olsection{Teorias consistentes com $\Th{Q}$ são indecidíveis}

O teorema a seguir afirma não apenas que $\Th{Q}$ é indecidível, mas
que, na verdade, toda teoria que não esteja em desacordo com $\Th{Q}$
é indecidível.

\begin{thm}
Seja $\Th{T}$ uma teoria qualquer na linguagem da aritmética que seja
consistente com $\Th{Q}$ (isto é, $\Th{T} \cup \Th{Q}$ é consistente).
Então $\Th{T}$ é indecidível.
\end{thm}

\begin{proof}
Lembre-se de que $\Th{Q}$ tem um conjunto finito de axiomas, $!Q_1$,
\dots, $!Q_8$. Podemos até substituí-los por um único axioma,
$!E = !Q_1 \land \dots \land !Q_8$.

Suponha que $\Th{T}$ seja uma teoria decidível consistente
com~$\Th{Q}$. Seja
\[
C = \Setabs{!A}{\Th{T} \Proves !E \lif !A}.
\]
Mostramos que $C$ seria uma separação computável de $\Th{Q}$ e
$\Th{\bar Q}$, uma contradição. Primeiro, se $!A$ pertence a $\Th{Q}$,
então $!A$ é demonstrável a partir dos axiomas de $\Th{Q}$; pelo
teorema da dedução, existe !!a{derivation} de $!E \lif !A$ na lógica de
primeira ordem. Logo, $!A$ pertence a~$C$.

Por outro lado, se $!A$ pertence a $\Th{\bar Q}$, então existe uma
prova de $!E \lif \lnot !A$ na lógica de primeira ordem. Se $\Th{T}$
também prova $!E \lif !A$, então $\Th{T}$ prova $\lnot !E$; nesse caso,
$\Th{T} \cup \Th{Q}$ é inconsistente. Mas estamos supondo que
$\Th{T} \cup \Th{Q}$ é consistente; portanto, $\Th{T}$ não prova
$!E \lif !A$ e, assim, $!A$ não pertence a~$C$.

Mostramos que, se $!A$ pertence a $\Th{Q}$, então pertence a $C$ e,
se $!A$ pertence a $\Th{\bar Q}$, então pertence a $\Complement{C}$.
Logo, $C$ é uma separação computável, que é a contradição procurada.
\end{proof}

Este teorema é muito poderoso. Por exemplo, implica:

\begin{cor}
  A lógica de primeira ordem para a linguagem da aritmética, isto é, o conjunto
  \[
    \Setabs{!A}{\text{$!A$ é demonstrável na lógica de primeira ordem}},
  \]
  é indecidível.
\end{cor}

\begin{proof}
A lógica de primeira ordem é o conjunto das consequências de
$\emptyset$, que é consistente com~$\Th{Q}$.
\end{proof}

\end{document}
